% Draft note for arXiv math.CO / cs.IT.
% Plain article class; no packages beyond amsmath, amssymb, amsthm, hyperref.
% Compile: pdflatex note.tex
\documentclass[11pt,a4paper]{article}

\usepackage{amsmath}
\usepackage{amssymb}
\usepackage{amsthm}
\usepackage{mathrsfs}
\usepackage[hidelinks]{hyperref}

\newtheorem{theorem}{Theorem}
\newtheorem{proposition}[theorem]{Proposition}
\newtheorem{corollary}[theorem]{Corollary}
\theoremstyle{definition}
\newtheorem{definition}[theorem]{Definition}
\theoremstyle{remark}
\newtheorem{remark}[theorem]{Remark}

\newcommand{\FF}{\mathbb{F}}
\newcommand{\HH}{\boldsymbol{H}}

\title{A binary $[50,40]_2$ code of covering radius $2$,\\
and a new seed for the $R=2$ iterative family}
\author{\footnotesize Grok Bot (4.6), Opus 5, Fable 5, GPT Sol 5.6\\
{\tiny with instructions from Stephen Wu}}
\date{16 August 2026}

\begin{document}
\maketitle

\begin{abstract}
We exhibit a binary linear $[50,40]_2$ code of covering radius $2$, showing
$\ell_2(10,2)\le 50$ and improving by one the length $51$ of the
Kaikkonen--Rosendahl code of 2003, which is the entry recorded for $r=10$ in
Table~5.1 of Davydov, Marcugini and Pambianco
(arXiv:2511.02542, November 2025). The parity-check matrix admits a
$(2,0)$-partition into $p(\HH)=10$ subsets, one fewer than the
computer-searched $p(\HH_{KR})=11$ of that paper. Since their infinite $R=2$
family is seeded by the $r=10$ entry, substituting this code moves the whole
iteration: Construction $\mathrm{QM}_2^2$ gives $[815,797]_22$ and
$[1631,1611]_22$ codes at $r=18$ and $r=20$, against the published $831$ and
$1663$, and iterating gives $n=51\cdot2^{r/2-5}-1$ for the reachable even $r$,
with asymptotic covering density $51^2/2^{11}\approx1.27002$ in place of
$26^2/2^{9}\approx1.32031$. All claims are exhaustively machine-checked from
the committed matrices; verification code accompanies the note.
\end{abstract}

\section{Introduction}

Let $\ell_2(r,R)$ denote the smallest length of a binary linear code with
codimension $r$ and covering radius $R$. If $\HH$ is an $r\times n$
parity-check matrix with column set $S\subset\FF_2^{\,r}$, then the code has
covering radius at most $2$ if and only if
\begin{equation}\label{eq:cover}
\{0\}\cup S\cup(S+S)=\FF_2^{\,r}.
\end{equation}
The covering density of an $[n,n-r]_2 2$ code is
$\mu=\bigl(1+n+\binom n2\bigr)/2^r$, and
$\bar\mu(2)=\liminf_{r\to\infty}\mu$ along a family.

Davydov, Marcugini and Pambianco~\cite{DMP2025} recently improved the upper
bounds on $\ell_2(r,2)$ for many $r$ and produced a new infinite family with
$\bar\mu(2)\approx1.32031$, superseding the value $\approx1.42383$ that had
stood since 1992. Their Table~5.1 records, for $r=10$, the length $51$ of the
Kaikkonen--Rosendahl code~\cite{KR2003}; that entry is not improved in their
paper.

The $r=10$ entry is load-bearing. Their family (\cite{DMP2025}, Theorem~5.7)
has lengths
\[
\Phi(r)=26\cdot2^{r/2-4}-1=2^{m}(51+1)-1,\qquad m=r/2-5,
\]
which is precisely the Kaikkonen--Rosendahl length $51$ pushed forward by
Construction $\mathrm{QM}_2^2$; their Theorem~5.7(i) states that the
$2$-partitions of $\HH_{KR}$ are ``very important for the next iterative
process''. A shorter seed admitting a small enough $2$-partition therefore
re-seeds the entire family.

\section{The code}

\begin{theorem}\label{thm:main}
Let $\HH$ be the $10\times50$ binary matrix listed in Appendix~\ref{app:H}.
Then $\HH$ has $\FF_2$-rank $10$, its $50$ columns are distinct and nonzero,
and its column set satisfies \eqref{eq:cover}. Hence $\HH$ is the parity-check
matrix of a binary $[50,40,3]_2 2$ code of covering radius exactly $2$, and
\[
\ell_2(10,2)\le 50 .
\]
Its covering density is $\mu=319/256=1.24609375$, against
$1327/1024\approx1.29590$ for the Kaikkonen--Rosendahl code.
\end{theorem}

\begin{proof}
Finite verification. Enumerating the empty sum, the $50$ singletons and all
$\binom{50}{2}=1225$ unordered pairs of columns and tabulating the resulting
$1276$ syndromes shows that all $1024$ elements of $\FF_2^{10}$ occur, with
multiplicity histogram
$\{1{:}859,\ 2{:}129,\ 4{:}24,\ 5{:}9,\ 6{:}3\}$.
Exactly $973$ syndromes are neither $0$ nor a column of $\HH$, so the covering
radius is not $1$; it is therefore exactly $2$. The rank and distinctness
claims are immediate from the listing. Minimum distance $3$ follows from the
$10$ linearly dependent triples reported in Proposition~\ref{prop:aux}.
\end{proof}

\begin{proposition}\label{prop:aux}
With $\HH$ as in Theorem~\ref{thm:main}:
\begin{enumerate}
\item[(i)] Deleting any single column of $\HH$ leaves at least $9$ elements of
$\FF_2^{10}$ uncovered; the best deletions, of the columns $381$, $479$ and
$927$, leave exactly $9$. Consequently the column set is a minimal
$1$-saturating set in $PG(9,2)$, equivalently a locally optimal covering code.
\item[(ii)] $\HH$ has minimum distance $d=3$ and exactly $10$ linearly
dependent triples of columns.
\end{enumerate}
\end{proposition}

\begin{proof}
Finite verification; for (i), $50$ recounts over all $2^{10}$ syndromes.
\end{proof}

\section{A $2$-partition with $p(\HH)=10$}

We use Definition~3.2 of~\cite{DMP2025}: a partition of the column set of
$\HH$ into nonempty subsets is a $(2,0)$-partition if every element of
$\FF_2^{\,r}$, including $0$, is a sum of at most two columns of $\HH$
belonging to distinct subsets, and $p(\HH)$ denotes the number of subsets.

\begin{theorem}\label{thm:part}
The partition $\mathscr P$ of the columns of $\HH$ into the ten subsets
\begin{align*}
&\{2,128,202,212,771,855,897,981\},\ \{86\},\ \{381,893,1003\},\ \{183,297\},\\
&\{1,65,247,256,320,438,502,734\},\ \{15,173,329,366,460,559,653,846\},\\
&\{4,16,391,491,742,754,869,881\},\ \{479,1004\},\ \{403\},\\
&\{32,341,373,576,608,777,789,821,927\}
\end{align*}
(columns written as integers, cf.\ Appendix~\ref{app:H}) is a
$(2,0)$-partition, so $p(\HH;\mathscr P)=10$. Moreover the columns
$491$, $734$, $821$ sum to the zero column and lie in three distinct subsets
of $\mathscr P$.
\end{theorem}

\begin{proof}
Finite verification. Of the $1024$ syndromes, one is zero and $50$ are columns
of $\HH$, both admissible by Definition~3.2; each of the remaining $973$ is
realised by a pair of columns from distinct subsets. For the last claim,
$491\oplus734\oplus821=0$ and these lie in the seventh, fifth and tenth subsets
as listed.
\end{proof}

By comparison, \cite{DMP2025} Theorem~5.2 obtains $p(\HH_{KR})=11$ by computer
search. The final assertion of Theorem~\ref{thm:part} is the analogue of their
Theorem~5.2(ii), which they require in order to apply Construction
$\mathrm{QM}_5^2$ at $r=28$.

\begin{remark}
The partition is tight in the following sense: $821$ of the $973$ syndromes
requiring a pair have a \emph{unique} such representation, and each of those
forces its two columns into distinct subsets.
\end{remark}

\section{Propagation}

Construction $\mathrm{QM}_2^2$ (\cite{DMP2025}, Theorem~4.1, equations (4.2)
and (4.4)) takes an $[n_0,n_0-r_0]_2 2$ code with a $2$-partition into
$p(\HH_0)$ subsets, an indicator set $\mathscr B=\FF_{2^m}$ subject to
$n_0\ge2^m\ge p(\HH_0)$, and $\boldsymbol D=\boldsymbol D_1(2)$, and returns an
$[n,n-r]_2 2$ code with
\[
n=2^m(n_0+1)-1,\qquad r=r_0+2m,\qquad p(\HH_C)\le2^{m+1}+1 .
\]
With $n_0=50$ and $p(\HH_0)=10$ the condition $n_0\ge2^m\ge p(\HH_0)$ holds
exactly for $m=4$ and $m=5$.

\begin{theorem}\label{thm:prop}
Applying $\mathrm{QM}_2^2$ to the code of Theorem~\ref{thm:main} with the
partition of Theorem~\ref{thm:part} yields
\begin{align*}
m=4:&\quad [815,797]_2 2,\ r=18,\ \mu=332521/262144\approx1.26847,\\
m=5:&\quad [1631,1611]_2 2,\ r=20,\ \mu=1330897/1048576\approx1.26924,
\end{align*}
so $\ell_2(18,2)\le815$ and $\ell_2(20,2)\le1631$. The published values are
$\Phi(18)=831$ and $\Phi(20)=1663$.
\end{theorem}

\begin{proof}
The two matrices are constructed explicitly and each is verified by complete
enumeration of $\FF_2^{18}$ and $\FF_2^{20}$ respectively: $262144/262144$ and
$1048576/1048576$ syndromes are covered, and both matrices have full rank.
\end{proof}

\begin{corollary}\label{cor:family}
Iterating $\mathrm{QM}_2^2$ from the seed of Theorem~\ref{thm:main}, a state
$(r,p)$ with length $n(r)=51\cdot2^{r/2-5}-1$ admits a step with parameter $m$
whenever $n(r)\ge2^m\ge p$, producing codimension $r+2m$ and
$p\le2^{m+1}+1$. The even $r\le64$ so reachable are
\[
10,\ 18,\ 20,\ 30,\ 32,\ 34,\ 36,\ 38,\ 40,\ 46,\ 48,\ 50,\ 52,\ 54,\ 56,\
58,\ 60,\ 62,\ 64,
\]
with $\ell_2(r,2)\le51\cdot2^{r/2-5}-1$ and
\[
\mu(r)=\frac{51^2}{2^{11}}-\frac{51}{2^{r/2+1}}+\frac{1}{2^{r}}
\ \nearrow\ \frac{51^2}{2^{11}}=\frac{2601}{2048}\approx1.27002 .
\]
Hence $\bar\mu(2)\le2601/2048$, and since Green's $f(2)$ is a
$\liminf$ over $r$, also $f(2)\le2601/2048$.
\end{corollary}

The even $r\in\{12,14,16,22,24,26,28,42,44\}$ are \emph{not} reachable by
$\mathrm{QM}_2^2$ from this seed. In~\cite{DMP2025} the analogous gaps at
$r=22,24,26$ are filled by Construction $\mathrm{QM}_3^2$ and $r=28$ by
$\mathrm{QM}_5^2$, with $r=42,44$ following from the $r=28$ code. The
ingredient that $\mathrm{QM}_5^2$ requires is present here, by the last
assertion of Theorem~\ref{thm:part}, but those constructions are not carried
out in this note and no claim is made for those $r$.

\section{What is not claimed}

The length $50$ is not shown to be optimal: the sphere-covering bound gives
only $\ell_2(10,2)\ge45$, since $1+n+\binom n2\ge1024$ fails at $n=44$.
Proposition~\ref{prop:aux}(i) shows only that this particular $50$-set has no
redundant column. An annealing run at $n=49$ terminated with $7$ uncovered
syndromes; that is search residue and not a lower bound.

Theorems~\ref{thm:main}, \ref{thm:part} and~\ref{thm:prop} are established by
complete enumeration. Corollary~\ref{cor:family} is not: beyond $m=4,5$ it
inherits the bound $p(\HH_C)\le2^{m+1}+1$ of~\cite{DMP2025} (4.4) rather than
exhibiting the partitions, exactly as the corresponding argument
in~\cite{DMP2025} does. Computing those partitions explicitly would strengthen
the corollary.

Table~5.1 of~\cite{DMP2025} is stated to be best ``as far as the authors
know''. We have not been able to search the ACCT proceedings exhaustively and
make no priority claim for $\ell_2(10,2)\le50$.

\section*{Verification}

The matrix of Appendix~\ref{app:H}, the partition of Theorem~\ref{thm:part},
and the two propagated matrices of Theorem~\ref{thm:prop} are accompanied by
two independent verifiers, in Python and in Rust, which read the matrices from
plain text and re-derive every number in this note by exhaustive enumeration:
no stored certificate is consulted and no sampling is used. The two use
opposite algorithms --- one enumerates pairs and marks the syndromes they hit,
the other sweeps the syndromes and searches for a partner column --- and their
outputs are compared automatically. The complete run takes about a second and
is deterministic.

The matrix was located by targeted simulated annealing seeded from
$\HH_{KR}$, run by an automated agent. The search procedure is not part of the
argument: Theorem~\ref{thm:main} is a finite decidable assertion about a fixed
matrix.

\appendix

\section{The matrix $\HH$}\label{app:H}

Rows $1$ to $10$ from top to bottom, columns left to right.

\begin{verbatim}
10010010011001010110111100110100100011110111111110
01010001001101000001001110111100011110001100001010
00110001011011000011111011101100111000111110101101
00010000010100010101010001110100110001001000101011
00001001001011000010110110101000010100110101101100
00000100011001010001110010011101001100010011100011
00000011000111001111110001111011011100001111100111
00000000111111000000001111111000111100000000011111
00000000000000111111111111111000000011111111111111
00000000000000000000000000000111111111111111111111
\end{verbatim}

\noindent The same columns as unsigned integers, with bit $i$ (least
significant first) equal to row $i+1$:

\begin{verbatim}
1 2 4 15 16 32 65 86 128 173 183 202 212 247 256 297 320 329 341
366 373 381 391 403 438 460 479 491 502 559 576 608 653 734 742
754 771 777 789 821 846 855 869 881 893 897 927 981 1003 1004
\end{verbatim}


\begin{thebibliography}{9}

\bibitem{DMP2025}
A.~A. Davydov, S.~Marcugini, F.~Pambianco,
\emph{New upper bounds for binary linear covering codes},
arXiv:2511.02542 [cs.IT], November 2025.

\bibitem{KR2003}
M.~Kaikkonen, C.~Rosendahl,
\emph{New covering codes from an ADS-like construction},
IEEE Trans.\ Inform.\ Theory \textbf{49} (2003), no.~7, 1809--1812.
\url{https://doi.org/10.1109/TIT.2003.813508}.

\bibitem{CHLL}
G.~Cohen, I.~Honkala, S.~Litsyn, A.~Lobstein,
\emph{Covering Codes}, North-Holland Mathematical Library \textbf{54},
Elsevier, Amsterdam, 1997.

\bibitem{Green}
B.~Green, \emph{100 open problems}, Problem 40.
Available at \url{https://people.maths.ox.ac.uk/greenbj/papers/open-problems.pdf}.

\end{thebibliography}

\end{document}
